\chapter{Parameters Characterizing a First-Order Phase Transition}
\label{ch:phase-transition-parameters}

A first-order cosmological phase transition is characterized by several
parameters that determine both its dynamics and its observable
signatures. Before defining them, we recall the relevant cosmological
background. The Robertson-Walker metric is
\begin{equation}
	ds^2 = dt^2 - a^2(t)\,d\mathbf{x}^2,
\end{equation}
with Hubble rate $H = \dot a/a$ satisfying the Friedmann equation
\begin{equation}
	H^2 = \frac{8\pi G}{3}\rho_{\text{total}},
	\qquad
	\rho_{\text{total}} = \frac{\pi^2}{30}g_* T^4 + \Delta\rho,
\end{equation}
where $g_*$ counts the relativistic degrees of freedom and $\Delta\rho$
is the contribution from vacuum energy. The temperature evolves as
\begin{equation}
	\frac{dT}{dt} = -HT,
	\label{eq:reference_temperature}
\end{equation}
which follows from entropy conservation in the expanding universe.

\section{Critical Temperature}

The critical temperature $T_c$ is the temperature at which the two
phases have equal free energy density, so the effective potential has
two degenerate minima. Above $T_c$ the symmetric phase is the global
minimum; below $T_c$ the broken phase becomes energetically preferred.
The universe does not transition instantaneously at $T_c$, however,
because the transition requires nucleation of bubbles of the true
vacuum, which is a tunneling process suppressed by $e^{-S_3/T}$.

% [FIGURE: effective potential at T > T_c, T = T_c, and T < T_c
%  showing the evolution of the two minima]

\section{Nucleation Temperature}

The number of bubbles nucleated per Hubble volume up to time $t$ is
\begin{equation}
	N(t) = \int_{-\infty}^{t}d\tilde t\,
	\Gamma(\tilde t)\,H^{-3}(\tilde t).
\end{equation}
Using \eqref{eq:reference_temperature} to change variables from $t$
to $T$,
\begin{equation}
	N(T) = \int_T^{\infty}\frac{d\tilde T}{\tilde T}\,
	\Gamma(\tilde T)\,H^{-4}(\tilde T).
\end{equation}
The nucleation temperature $T_n$ is defined by $N(T_n) = 1$, the
condition that on average one bubble has nucleated per Hubble volume,
\begin{equation}
	\Gamma(T_n)\,H^{-4}(T_n) \sim 1.
\end{equation}
For $O(3)$-symmetric instantons at finite temperature,
$\Gamma(T) \sim T^4 e^{-S_3/T}$, which gives the approximate condition
\begin{equation}
	\frac{S_3(T_n)}{T_n} \approx 145
	- 4\ln\!\left(\frac{T_n}{100\,\text{GeV}}\right).
\end{equation}
This determines $T_n$ once $S_3(T)$ is known from the bounce
calculation.

\section{The Parameter \texorpdfstring{$\tilde\beta$}{beta}}

The parameter $\tilde\beta$ characterizes how fast the nucleation rate
grows near $T_n$,
\begin{equation}
	\tilde\beta = \frac{1}{\Gamma}\frac{d\Gamma}{dt}
	\bigg|_{N=1},
\end{equation}
so that $\tilde\beta^{-1}$ gives the characteristic duration of the
phase transition. Using \eqref{eq:reference_temperature},
\begin{equation}
	\tilde\beta = -HT\frac{d\ln\Gamma}{dT}\bigg|_{T=T_n}.
\end{equation}
Writing $\Gamma = A(T)e^{-S_E(T)}$ and neglecting the slowly varying
prefactor $A(T)$,
\begin{equation}
	\boxed{\frac{\tilde\beta}{H}\bigg|_{T_n}
		= T_n\frac{dS_E}{dT}\bigg|_{T=T_n}.}
\end{equation}
This is a dimensionless number. When $\tilde\beta/H \gg 1$ the phase
transition completes in a time much shorter than a Hubble time, so
the universe can be treated as approximately static during the
transition and the temperature does not change appreciably.

\section{The Parameter \texorpdfstring{$\alpha$}{alpha}}

The latent heat released during the transition is the energy density
difference between the two phases,
\begin{equation}
	\Delta\rho = \left(\Delta V - T\frac{\partial\Delta V}{\partial T}
	\right)_{T=T_n}
	\approx \Delta V\big|_{T=T_n},
\end{equation}
where $\Delta V$ is the free energy difference and the derivative term
is subleading. The parameter $\alpha$ is defined as the ratio of this
energy density to the radiation energy density,
\begin{equation}
	\alpha = \frac{\Delta\rho}{\rho_{\text{rad}}}\bigg|_{T=T_n},
	\label{eq:alpha}
\end{equation}
where $\rho_{\text{rad}} = \pi^2 g_* T_n^4/30$. The parameter $\alpha$
measures the strength of the transition: $\alpha \ll 1$ corresponds to
a weak transition in which the latent heat is a small perturbation on
the radiation bath, while $\alpha \gg 1$ corresponds to a strongly
supercooled transition.

% [FIGURE: effective potential showing Delta V, the energy gap between
%  the two phases at T = T_n]

\section{Reheating Temperature}

In the supercooled case $\alpha \gg 1$, the latent heat released by
the transition dominates over the radiation energy density. The plasma
is reheated to a temperature $T_{\text{rh}} \gg T_n$ determined by
energy conservation,
\begin{equation}
	\left(\Delta\rho + \rho_{\text{rad}}\right)_{T=T_n}
	= \frac{\pi^2}{30}g_* T_{\text{rh}}^4.
\end{equation}
This reheating can have important consequences for baryogenesis and
other processes that are sensitive to the thermal history of the
universe.

\section{Bubble Wall Velocity}

The bubble wall velocity $v_w$ is determined by the balance between
the pressure difference driving the wall outward and the friction
from particles in the plasma acquiring mass as they cross the wall.
It is intrinsically a non-equilibrium quantity and difficult to
calculate from first principles. Three regimes are distinguished.

In the non-relativistic deflagration regime, $v_w \ll 1$, the wall
moves slowly and the plasma in front of the wall has time to
equilibrate. In the relativistic but non-runaway regime, $v_w \to 1$
but the Lorentz factor $\gamma = (1-v_w^2)^{-1/2}$ approaches a
finite constant, as the friction force grows to balance the driving
pressure. In the runaway regime, $v_w \approx 1$ and $\gamma$
grows without bound throughout the transition.

In theories with a Higgs boson, the Higgs acquires a mass as it
crosses the bubble wall, going from zero mass in the false vacuum to
a finite mass in the true vacuum. The resulting friction force is
generically large enough to prevent the runaway regime from being
reached, so the bubbles at most attain relativistic but finite
velocities (Bodeker and Moore, arXiv:1703.08215).

The parameters $\alpha$, $\tilde\beta$, and $v_w$ together determine
the spectrum of gravitational waves produced by the phase transition,
through the three production mechanisms of bubble wall collisions,
sound waves in the plasma, and magnetohydrodynamic turbulence.
